\chapter{Release}
\ChapterLabel{Release}
\section{Chapter Responsibility}
\label{Section:ReleaseResponsibility}
Identify the exact deliverable and establish its readiness for authorized release.

\noindent\textbf{Expected result:} A controlled release package presented to the release authority.

\section{How to Use This Chapter}
Identify the release and define contents, approvals, evidence, known limitations, deployment instructions, rollback or recovery approach, and support ownership.
\section{Week 2 Requirements Baseline for Technical Review}
\label{Section:WeekTwoReviewPackage}
\textbf{Package:} Embedded ADC System, document draft 0.4 (Week 2).
\textbf{Working branch:} \texttt{week2}; \textbf{review target:}
\texttt{main}. \textbf{Prepared:} November 12, 2026.
\textbf{Owner:} Jack Marino. The exact review revision is the pull request's
head commit. This package is a proposed requirements baseline presented for
technical review; neither requirements approval nor release authorization
has been granted.

\paragraph{Review contents.}
The package contains the authored LaTeX, editable TikZ sources, bibliography,
and complete-document CI workflow. It includes six current stories in
\ChapterRef{UserStories}, 40 candidate requirements in
\hyperref[Section:WeekTwoCandidates]{the requirements chapter}, their candidate
use cases and planned checks, 40
\hyperref[Table:WeekTwoVerification]{verification entries}, six
\hyperref[Table:ValidationRegister]{validation entries}, and 46
\hyperref[Section:PlannedEvidence]{evidence collection plans}. The
\hyperref[Table:TraceabilityMatrix]{trace matrix} links those artifacts and
retains the explicit portability coverage gap for \UserStoryRef{US-03}.
The package also carries the ADC scope/context and build support needed for
the complete architecture book on the review target.

\paragraph{Technical review focus.}
Review requirement necessity, singularity, units, feasibility, scenario
boundaries, hardware capture and clock assumptions, finite load/buffering
envelopes, and the adequacy of planned acceptance observations. Review the
candidate local record contract separately from
\hyperref[Section:ProposedOMSBoundary]{the OMS decision record}.
\RequirementRef{REQ-OMS-01} remains Proposed; applicability, standard version,
CAL profile, UCI message, status mapping, responsible gateway, and end-to-end
contract are TBD. Jack Marino owns unresolved decisions; decision dates remain
TBD as agreed. The lecture's 10 ms, 1 ms, and 100 ms allocations are teaching
values pending project approval, not product commitments.

\paragraph{Build and evidence boundary.}
Build the complete book using the platform script in
\ChapterRef{DocumentGuide}, including bibliography, glossary, index, and final
reference-resolution passes. The CI workflow also runs Latexmk, checks for
unresolved references/citations, and retains the compiled PDF, logs, text,
layout warnings, and review-page renders. The latest successful run for the
review head is document-build evidence only. Engineering verification and
validation remain Pending; planned evidence remains Not collected.
\textbf{Technical review approval: TBD; owner: Jack Marino; decision date: TBD.}
Record reviewer decisions and resulting changes before approving a baseline
or authorizing release.

\section{Retained Template Checklist}
Confirm the version and artifact IDs are current; all required traceability is complete; required tests, verification, and validation pass or have approved waivers; release notes and known limitations are published; and the release authority approves the exact package.
